\subsection{Format Preserving Encryption}\label{sec:symmetric_cryptography:format_preserving}

Most symmetric ciphers (\Cref{sec:symmetric_cryptography:ciphers}) do not necessarily preserve the format of the plaintext --- intended as the set of allowed characters and overall structure --- when encrypting. For instance, encrypting a credit card number usually produces a (pseudo-)random binary string rather than a sequence of digits only. The loss of format can pose challenges to existing (especially legacy) information systems that may struggle to accommodate changes in data format. To solve this problem, \textbf{\gls{fpe}} --- a category of symmetric ciphers which preserve plaintexts' format --- was developed. 

Early \gls{fpe} ciphers were built on standard block ciphers (\Cref{sec:cryptography_basics:ciphers}) where plaintext and (keystreams derived from) keys were defined over a restricted set of characters called \textbf{alphabet} \cite{nist_guidelines_1981}. Then, Black and Rogaway proposed the ``cycle‐following'' method \cite{black_ciphers_2002}, which consists of repeating the encryption process iteratively until a ciphertext is produced that conforms to the desired format. The cycle-following method was refined over time and --- at the time of writing --- the state-of-the-art includes the FF1 and FF3-1 algorithms, recommended by \gls{nist} \cite{dworkin_recommendation_2025}; both FF1 and FF3 are instantiations of the FFX framework \cite{bellare_ffx_2010}. Another \gls{fpe} standard is ANSI ASC X9.124 by major organizations including American Express, IBM, MasterCard, \gls{nist}, Thales, and VISA. \cite{ansi_fpe_2023}; ANSI ASC X9.124 specifies three \gls{fpe} schemes derived from FF1 and FF3, with predefined secure parameters and configurations.

\remarkBlock{Symmetric encryption vs. \gls{fpe}}{\gls{fpe} is generally regarded as less robust and less efficient than non-\gls{fpe} symmetric ciphers, as evidenced by the numerous vulnerabilities that have been identified \cite{bellare_message_2016,shacham_curse_2018}.}
